% Sources: ultrametric-regularisation-and-loss/body.tex and unification.tex
% Thesis commit: 2c6418bcf9643fc6e039237f0f59ace14b2557fc

\begin{thm}[Discrete Metric Regularised Regression]
\label{thm:discrete-regularised}
For regularised linear regression with the discrete metric and regularisation strength $\lambda \geq 0$, an optimal hyperplane $(a_1, \ldots, a_n, b)$ has exactly $t^*$ non-zero coefficients among $a_1, \ldots, a_n$, where
\begin{equation}
    t^* \in \argmax_{t \in \{0, \ldots, n\}} \left( K(t) - \lambda t \right)
\end{equation}
and the corresponding minimum loss is $m - K(t^*) + \lambda t^*$.
\end{thm}

\begin{corollary}[Reduction of regularised to unregularised search]
\label{cor:discrete-algorithm}
Given the values $K(0), K(1), \ldots, K(n)$, together with a witness
hyperplane from $B(t)$ for each $t$, an optimal hyperplane for regularisation
strength $\lambda$ can be found by:
\begin{enumerate}
    \item Computing the $n+1$ candidate gains: $\{K(0), K(1)-\lambda, K(2)-2\lambda, \ldots, K(n)-n\lambda\}$
    \item Selecting any $t^*$ that maximises $K(t^*) - \lambda t^*$
    \item Returning the witness hyperplane for $B(t^*)$
\end{enumerate}
\end{corollary}

\begin{thm}[Additive contact theorem]
\label{thm:additive-contact}
Let \(\{(X_i,y_i)\}_{i=1}^{m}\subseteq \mathbb{Q}^n\times \mathbb{Q}\) satisfy
the hypotheses of Theorem~\ref{core-theorem}, and write
\[
r_i(\beta)=\beta\cdot (X_i,1)-y_i.
\]
Let \(\phi:[0,\infty)\to\mathbb{R}\) be nondecreasing and satisfy
\(\phi(0)<\phi(t)\) for every \(t>0\). Define
\[
L_{\phi}(\beta)=\sum_{i=1}^{m}\phi\bigl(|r_i(\beta)|_p\bigr).
\]
Then any minimiser of \(L_{\phi}\) passes through at least \(n+1\) data points,
unless the data are degenerate in the sense of Theorem~\ref{core-theorem}.
\end{thm}

\begin{corollary}[Count loss and the additive \(q\)-family]
\label{cor:additive-contact-special-cases}
Under the same hypotheses, every minimiser of
\[
L_{\mathrm{count}}(\beta)=\sum_{i=1}^{m}\mathbf{1}_{\{r_i(\beta)\neq 0\}}
\]
and every minimiser of
\[
\mathcal{L}_q(\beta)=\sum_{i=1}^{m} q^{-v_p(r_i(\beta))} \qquad (q>1)
\]
passes through at least \(n+1\) data points, unless the data are degenerate.
\end{corollary}

\begin{thm}[When \(q>m\), additive loss is lexicographic]
\label{thm:q-lexicographic}
Fix a regression instance with \(m\) residuals, and let
\[
h_t(\beta):=\#\{\,i : v_p(r_i(\beta))=t\,\}
\qquad (t\in\mathbb{Z}).
\]
Let \(q>m\). Suppose \(\beta\) and \(\beta'\) have different valuation
histograms, and let \(t_0\) be the smallest integer for which
\[
h_{t_0}(\beta)\neq h_{t_0}(\beta').
\]
If
\[
h_{t_0}(\beta)<h_{t_0}(\beta'),
\]
then
\[
\mathcal{L}_q(\beta)<\mathcal{L}_q(\beta').
\]
Equivalently, for \(q>m\), minimising \(\mathcal{L}_q\) is the same as
lexicographically minimising the residual-valuation histogram from the worst
valuation upward.
\end{thm}

\begin{thm}[Max-loss contact existence]
\label{thm:max-contact-existence}
Consider the sup-norm objective
\[
\mathcal{L}_{\max}(\beta)=\max_{i}|r_i(\beta)|_p.
\]
Assume the data satisfy the hypotheses of Theorem~\ref{core-theorem} and are
non-degenerate. Then \(\mathcal{L}_{\max}\) attains its minimum, and at least one
minimiser passes through at least \(n+1\) data points. (Unlike the additive
case, not every minimiser need have \(n+1\) contacts.)
\end{thm}
